%=============================================================================!
% 14.9  WHAT THE TRAJECTORY LOOKS LIKE                                        !
%=============================================================================!
\section{What the trajectory looks like}
\label{sec:pr-trajectory}

A run at constant pressure has two more things to read than one at fixed
volume: the volume itself, and a pressure whose single values mean little.
This section shows a healthy run of the liquid and six faults that
pressures and barostats leave, each with the check that catches it.

\subsection*{A healthy run}

\Cref{fig:pr-healthy} follows the liquid for \SI{300}{\pico\second} under
stochastic cell rescaling at \SI{0.1}{\giga\pascal} with $\tau_{\mathrm P}
= \SI{1}{\pico\second}$ and CSVR at \SI{135}{\kelvin}, started from the box
of Chapter~13 at \SI{12577}{\angstrom\cubed}. After the first
\SI{10}{\pico\second} the volume averages \SI{12312}{\angstrom\cubed}, a
density of $0.8172\sigma^{-3}$, and spreads by \SI{190}{\angstrom\cubed},
against the \SI{183}{\angstrom\cubed} of \cref{eq:pr-volume}; it lies within
the predicted \SI{183}{\angstrom\cubed} of its mean for 0.667 of the time,
close to the 0.683 of a Gaussian (\cref{sec:sm-probability}). The instantaneous pressure scatters by \SI{0.0175}{\giga\pascal}
about its mean of \SI{0.0988}{\giga\pascal}: for 256 atoms a single value of
$P$ says little, and only its average is compared with the target. The
sum $K + U + P_0V$ wanders over \SI{9.50}{\milli\electronvolt} per atom as
the thermostat and the barostat exchange energy with the liquid, while the
effective energy, that sum less the heat and the work, stays within
\SI{0.027}{\milli\electronvolt} per atom and drifts by
\SI{-6.4e-6}{\electronvolt\per\pico\second} for the whole liquid. The kinetic temperature
averages \SI{134.69}{\kelvin}.

\begin{figure}[htb]
\centering
\includegraphics{ch14_pressure/healthy}
\caption{A healthy run at constant pressure: the liquid under stochastic
cell rescaling at \SI{0.1}{\giga\pascal} and CSVR at \SI{135}{\kelvin}, for
\SI{300}{\pico\second}. (a) The volume every \SI{100}{\femto\second}, with
its mean and the spread $\sqrt{\kB T\ens{V}\kappa_T}$ either side (dotted).
(b) The instantaneous pressure (thin) and its running mean (thick),
against the target (dotted). (c) $K + U + P_0V$ per atom and the effective
energy, that less the heat and the barostat's work, from their starting
values.}
\label{fig:pr-healthy}
\end{figure}

\subsection*{Six faults}

\Cref{fig:pr-faults} shows six faults and \cref{tab:pr-faults} the check
that catches each.

\begin{figure}[p]
\centering
\includegraphics{ch14_pressure/faults}
\caption{Pressure and barostat faults (teal) against healthy runs (grey
dashed). (a) The pressure of the liquid at fixed volume from
$\sum_i\vb{r}_i\cdot\vb{F}_i$ with wrapped positions, against the pair
form. (b) The volume under cell rescaling given the virial's pressure
without the kinetic part. (c) The density of the volume under
Berendsen's barostat, against cell rescaling. (d) The volume under a
piston with $\tau_{\mathrm P} = \SI{10}{\pico\second}$, against
\SI{1}{\pico\second}. (e) The layered solid with guests: the pressure in
its plane and across it under isotropic coupling, against
semi-isotropic coupling (grey, zero). (f) The running mean of the
pressure under Berendsen's barostat given $\kappa$ in bar$^{-1}$ where
GPa$^{-1}$ was meant, against the right $\kappa$, with the target
(dotted).}
\label{fig:pr-faults}
\end{figure}

\begin{table}[htb]
\centering
\caption{The faults of \cref{fig:pr-faults}, the check that catches each,
and its value in the healthy and the broken run.}
\label{tab:pr-faults}
\small
\begin{tabular}{@{}>{\raggedright\arraybackslash}p{0.27\linewidth}
   >{\raggedright\arraybackslash}p{0.33\linewidth}ll@{}}
\toprule
fault & check & healthy & broken\\
\midrule
(a) virial from wrapped positions & $P$ before and after moving the
   origin, GPa &
   0.0877, 0.0877 & 0.0324, 0.0354\\
(b) kinetic part dropped & $\ens{P}$ with $2K$ included, GPa & 0.0988 &
   0.1411\\
(c) volume too narrow & $\kappa_T$ from the fluctuations, GPa$^{-1}$ &
   1.50 & 0.64\\
(d) piston too heavy & swings of $V$ in \SI{90}{\pico\second} & 103 & 10\\
(e) isotropic coupling of layers & $\frac12(\mathsf{P}_{xx} +
   \mathsf{P}_{yy})$, $\mathsf{P}_{zz}$, GPa & 0.000, 0.000 & $-0.141$,
   0.282\\
(f) $\kappa$ in the wrong unit & $\ens{P}$ after the first
   \SI{50}{\pico\second}, GPa & 0.1001 & 0.0871\\
\bottomrule
\end{tabular}
\end{table}

\emph{(a) The virial from wrapped positions.} As \cref{sec:pr-periodic}
showed, $\sum_i\vb{r}_i\cdot\vb{F}_i$ with wrapped positions gives
\SI{0.0324}{\giga\pascal} for a liquid whose pressure is
\SI{0.0877}{\giga\pascal}, and a different value when the origin moves.
The check is to move the origin: a correct pressure does not change.

\emph{(b) The kinetic part dropped.} A barostat handed only $\mathcal{W}/3V$,
as happens when ASE's \texttt{get\_stress} is called without
\texttt{include\_ideal\_gas=True}, holds that part at $P_0$ and lets the
true pressure rise by $N\kB T/V$, here \SI{0.0408}{\giga\pascal}: the
liquid shrinks by 4.9\%, to \SI{11706}{\angstrom\cubed} from
\SI{12312}{\angstrom\cubed}, and its full pressure reads
\SI{0.1411}{\giga\pascal}. The check is the mean of $(2K +
\mathcal{W})/3V$ against $P_0$.

\emph{(c) A volume too narrow.} Berendsen's barostat holds the mean
pressure, so nothing in one run looks wrong; the check is $\kappa_T$ from
the volume's fluctuations, \cref{eq:pr-volume}, against an independent
value: \SI{0.64}{\per\giga\pascal} under Berendsen's barostat,
\SI{1.50}{\per\giga\pascal} under cell rescaling, and
\SI{1.46}{\per\giga\pascal} from the runs at fixed volume.

\emph{(d) A piston too heavy.} With $\tau_{\mathrm P} =
\SI{10}{\pico\second}$ the volume swings with a period of
\SI{8.7}{\pico\second}, 10 times in \SI{90}{\pico\second}, so the run's
mean volume rests on ten swings; with \SI{1}{\pico\second} it swings 103
times. The check is the period of the volume against the length of the
run.

\emph{(e) Isotropic coupling of a layered solid.} The mean pressure reads
the target while the solid is pulled along its sheets and squeezed across
them, $-0.141$ and \SI{+0.282}{\giga\pascal}. The check is the diagonal
of the pressure tensor, component by component.

\emph{(f) The compressibility in the wrong unit.} The barostat's $\kappa$
for the liquid, \SI{2.5}{\per\giga\pascal}, is $2.5\times10^{-4}$ per bar; entered
as $2.5\times10^{-4}$ where gigapascals were meant, it makes Berendsen's
barostat $10^4$ times slower, a relaxation time of
\SI{5.8}{\nano\second} instead of \SI{0.58}{\pico\second}. The volume
stays where it started, \SI{12577}{\angstrom\cubed}, and the pressure
stays at \SI{0.0871}{\giga\pascal} against the target 0.1. The check is
the running mean of the pressure against the target, and the volume's
relaxation against $\tau_{\mathrm P}$.

\begin{practice}
A run at constant pressure is read through the running mean of its
pressure, computed with the kinetic part and from pair separations, which
should settle at $P_0$; through the spread of its volume, which should be
$\sqrt{\kB T\ens{V}\kappa_T}$, with $\kappa_T$ checked at fixed volumes;
through the diagonal of its pressure tensor, which for a solid should
match $P_0$ in every component the coupling lets relax; and through the
period of its volume, or its relaxation time when the barostat has no
piston, against the length of the run. Berendsen's barostat serves to
bring a system to its pressure; stochastic cell rescaling or a piston
with chains then samples it, and $\tau_{\mathrm P} = \SI{1}{\pico\second}$
sampled the liquid's volume correctly here, with 103 swings of the piston
in \SI{90}{\pico\second}; a layered solid needs semi-isotropic coupling, or
anisotropic coupling when the directions of its plane differ.
\end{practice}

\begin{keyidea}
In a healthy run at constant pressure the instantaneous pressure
scatters widely, its running mean settles at the target, the volume
spreads by $\sqrt{\kB T\ens{V}\kappa_T}$, and the effective energy holds
still. A virial from wrapped positions, a pressure without its kinetic
part, Berendsen's narrow volume, a piston too heavy, isotropic coupling of
a solid that wants to change shape, and a compressibility in the wrong
unit each leave a mark that a single check catches.
\end{keyidea}

\notebook{14}{break the run at constant pressure in each of the six ways and read the checks}

\begin{selfcheck}
For a liquid of 500 atoms at \SI{300}{\kelvin} in
\SI{20000}{\angstrom\cubed}, minus the mean of the diagonal of ASE's
\texttt{get\_stress()}, called with its defaults, averages
\SI{-0.04}{\giga\pascal}. What is the liquid's pressure?
\end{selfcheck}
